\setcounter{section}{0}

{\Large\bfseries\filcenter
Facilities, Equipment, \& Other Resources
\\\rule{\textwidth}{0.4pt}}


\section{Facilities}




The PIs' office and his lab, \textbf{AI and Secure Computing (AISeC)}, are located next to each other in the Department of Computer Science at Loyola University Chicago. 
The AISeC lab has three desks and a larger one that can seat at least 10 people.   %The Cudahy Science building contains a 3D printing room, which all students will have access to.
The lab is equipped with multiple GPU-enabled workstations and desktops, and a large monitor for group meetings and presentations.
This work is co-hosted with the  \textbf{Software and Systems Laboratory (SSL)}, which has additional workstations, desks, and robotic equipment.
All infrastructure is maintained by Mr.~Miao Ye, our full-time systems and lab manager.



\section{Major Equipment for Project's Tasks}
In addtion to the GPU-enabled workstations in the AISeC and SSL labs (totaling +8 GPUs, +2TB RAM, +200 CPU Cores), the PI has access and support to the CS department's computing infrastructure, which is housed at the Research Data Center.
The department maintains four 42U racks of servers to address general-purpose computing, virtualization, and high-performance computing needs.
Virtualization is addressed through several racks dedicated to VMWare vSphere and Windows Hyper-V. 
Multiple racks are dedicated to shared storage and backup (via replication) with nearly 2 PB of disk space and growing.`'
High-performance computing is supported via a cluster (acquired in 2018) with 16 quad-core nodes, including support for GPU computing. The department also acquired multiple multi-GPU systems with Lambda stack support, which are used specifically for teaching and research needs in machine learning.
This equipment ensures that there are adequate resources for all tasks that will be undertaken for the proposed work. 

%The Department of Computer Science runs multiple ESXi and Windows servers that Virtual Machines (VMs) can be spawned on demand. These servers could be used to run any VM for classes and to participate in cybersecurity competitions. Other research activities, such as analyzing malware and running long-term experiments, can also be run on these VMs. A GPU cluster is also available for use by students. 


\section{Loyola Research Data Center}

Loyola University Chicago's Research Data Center (RDC) is a 1,000-square-foot facility dedicated to supporting research and funded projects and provides a secure home for the computational clusters and related equipment used by our research community.
The RDC, opened in 2010, provides a high-availability computing environment for research projects. This facility is equipped with power protection, including an uninterruptible power supply and a backup generator. 
Multiple computer room air conditioning (CRAC) units provide redundant cooling for the space, and a structured cabling design allows high-speed network connectivity. In addition to fire protection, additional safety and security elements for the RDC include keycard access, camera surveillance, and environmental monitoring.
Sized to accommodate moderate growth, several research initiatives are currently taking advantage of the space, which currently houses three research clusters and more than 100 nodes. 
Furthermore, collaborative research efforts with other participating institutions and organizations have full access and connectivity to the Internet via the Metropolitan Research \& Education Network (MREN) to accommodate high-bandwidth applications, data transmissions, and computational requirements.
A steering committee composed of senior administrators, faculty, and ITS professionals is responsible for reviewing, evaluating, and recommending strategies, plans, and policies governing the use of RDC resources. Loyola's RDC is managed by Information Technology Services (ITS) in partnership with the University's Facilities Department.




\section{Resources for Education and Outreach Activities}
Loyola University Chicago has a strong commitment to supporting research and education activities. The PI will have access to venues for hosting workshops, seminars, and outreach activities. The university has multiple large venues for hosting workshops and events, which the PI can and has already utilized in the past to host seminars. 
The university also supports undergraduate research via multiple initiatives, including the LUROP program, which allows students to engage in funded research with no additional grant cost.

% The PI will also have access to the Loyola Center for Experiential Learning (CEL) to support undergraduate research activities. CEL provides funding and support for undergraduate research assistants, including training and professional development.

% \section{Other Resources}
% The PI's lab contain a few desktops, laptops, and tablets that students can use for any activities related to this proposal. There are also ``docking'' monitors where students can plug in their laptop via USB-C for display and power.

% Loyola is a Microsoft customer and has access to the Microsoft Azure cloud platform which participants can use for the proposed hands-on activities.



